\section{Experiments}\label{sec:experiments}

We evaluate the proposed correspondence-generation pipeline in two complementary settings. First, we assess whether the generated correspondences can improve real-image semantic correspondence models when used as synthetic pretraining data (\autoref{spair71k_eval}). Second, we directly evaluate the quality of the generated 3D correspondences against human-annotated 3D ground truth (\autoref{annotation_quality}).
All experiments were conducted on a single NVIDIA RTX A6000 GPU.

\subsection{Experimental Data}
\label{sec:exp_data}

Our method assumes collections of rigid 3D objects that are approximately rotation-aligned, providing a weak category-level geometric prior. In this work, we use two complementary 3D sources, 
\textbf{KeypointNet}~\cite{you2020} and \textbf{Objaverse-OA}~\cite{orient_anything}, 
to validate the proposed pipeline. 
%
\textbf{KeypointNet} provides sparse human-annotated semantic keypoints on approximately aligned 3D object models. It contains relatively few rigid categories, but many objects per category, making it suitable for evaluating how well the correspondences generated by our pipeline agree with human semantic annotations. These annotations are used only for evaluation, not for training or correspondence generation.
%
\textbf{Objaverse-OA} provides substantially broader category and shape diversity, containing over 100 rigid categories after curation, but without correspondence annotations. \autoref{sup:qualitative}  shows example correspondences generated by our method on several Objaverse-OA categories beyond those covered by existing correspondence datasets. For the downstream SPair-71k experiments, we select rigid categories corresponding to the evaluated SPair-71k \cite{min2019spair71k} classes, exclude non-rigid, articulated, and strongly rotationally symmetric objects, and apply our pipeline to generate the synthetic correspondences used for pretraining.
%
Thus, the two sources play complementary roles: KeypointNet is used to evaluate annotation quality, while Objaverse-OA is used to test whether the generated correspondences provide useful synthetic supervision for downstream 2D semantic correspondence models.
\subimport{../tables}{spair_table}

\subsection{Results on SPair-71k}
\label{spair71k_eval}
We evaluate the synthetic 2D correspondence supervision generated by our pipeline on SPair-71k~\cite{min2019spair71k}, a standard benchmark for cross-instance semantic correspondence.
Concretely, we generate 12K synthetic 2D training pairs per category from diverse 3D object instances sampled from Objaverse-OA~\cite{orient_anything}. To preserve category and shape diversity in the synthetic pretraining data, we avoid pair-level semantic filtering in this experiment. We focus on 9 rigid categories with well-defined 3D semantic-geometric structure: aeroplane; bicycle; boat; bus; car; chair; motorbike; train; tvmonitor. We report reproduced results for recent methods under our 9-category protocol, including GeoAware-SC~\cite{zhang2024tellingLeftFromRight}, Jamais Vu~\cite{mariotti2025jamaisvu}, and MARCO~\cite{cuttano2026marco}, together with reported SPair-71k results for earlier methods.

For the state-of-the-art GeoAware-SC and MARCO, we also evaluate the effect of our generated data in three training regimes: training on SPair-71k only; joint training with our projected correspondences; pretraining on our generated correspondences followed by fine-tuning on SPair-71k. To preserve category and shape diversity during pretraining, we avoid pair-level filtering in this setting.

Since our data is generated from synthetic 3D objects, a synthetic-to-real domain gap remains, and naive joint training gives mixed results. In contrast, the standard pretrain-and-finetune regime consistently improves both evaluated models. GeoAware-SC~\cite{zhang2024tellingLeftFromRight} improves from 79.36 to 83.73 mean PCK (+4.37), and MARCO~\cite{cuttano2026marco} improves from 86.39 to 89.28 (+2.89). These gains are obtained without any architectural modification, highlighting that even approximate semantic-geometric correspondences provide a useful geometric-semantic training signal that is easier to obtain than manual annotations. Additional qualitative examples in \autoref{fig:qual_marco_main} further illustrate the improved semantic-geometric consistency obtained after pretraining on our generated data.


\subsection{Evaluation of Annotation Quality}
\label{annotation_quality}

We evaluate the accuracy of our generated 3D correspondences using KeypointNet~\cite{you2020}, which provides human-annotated semantic keypoints on approximately rotation-aligned 3D object models. This evaluation directly measures how well our predicted correspondences agree with human-defined semantic matches. For each category, we sample pairs of models and treat the keypoints of one model as queries, predicting their corresponding locations on the other model. We then compare the predicted correspondences to the ground-truth annotations using PCK.

Our main results use both correspondence-level semantic pruning and pair-level semantic filtering. Correspondence-level pruning removes low-confidence matches within a pair, while pair-level filtering removes object pairs with insufficient shared semantic support. We report results on three subsets of 100 models per category, and present the average performance per category in \autoref{tab:keypointnet_table}. Overall, our method achieves 93\% PCK@0.10, 97\% PCK@0.15, and 99\% PCK@0.20, indicating that although the generated correspondences are approximate rather than exact point-level matches, most predictions remain very close to the human-annotated semantic targets.

\subimport{./}{keypointnet_exps}

\subsection{Ablations and Hyperparameter Sensitivity}
\label{sec:ablation}

We summarize the main findings of our ablation studies here and defer detailed quantitative analyses, per-category tables, and additional sensitivity experiments to \autoref{sup:ablation_study} and \autoref{Sup:per_label}.

\noindent\textbf{Hyperparameters.}
Performance is stable across a reasonable range of the CPA smoothness weight \(\lambda_{\mathrm{smooth}}\) and semantic threshold \(\tau_{\mathrm{sem}}\). As detailed in \autoref{sup:ablation_semantic_filtering}, \(\tau_{\mathrm{sem}}\) controls an accuracy--coverage trade-off, and \(\tau_{\mathrm{sem}}=0.85\) provides a practical balance between retaining reliable correspondences and preserving sufficient coverage. For the CPA smoothness weight, we observe only minor variation in KeypointNet \cite{you2020} PCK across tested values, likely because the sparse annotations are mostly concentrated on salient corners and extremities. Qualitatively, however, \(\lambda_{\mathrm{smooth}}=0.3\) better preserves the source object structure while still allowing sufficient local deformation to align with the target, producing more coherent dense correspondences in regions not captured by the sparse KeypointNet annotations. We therefore use \(\lambda_{\mathrm{smooth}}=0.3\) and \(\tau_{\mathrm{sem}}=0.85\) in all experiments.

\noindent\textbf{Registration model.}
We compare continuous piecewise-affine (CPA) registration with a global affine transformation and thin-plate splines (TPS); see \autoref{sup:ablation_registration} for details. Affine performs substantially worse at all PCK thresholds, while TPS achieves accuracy closer to CPA. CPA nonetheless provides the best overall performance and is approximately \(3\times\) faster in our setup, likely because its tetrahedral piecewise-affine parameterization is local, whereas TPS relies on a global radial-basis deformation.

\noindent\textbf{Semantic filtering.}
Semantic filtering consistently improves correspondence quality (see \autoref{sup:ablation_semantic_filtering}). Without correspondence-level or pair-level filtering, geometric registration alone achieves over \(84\%\) PCK@0.10. Adding correspondence-level semantic pruning improves performance to about \(91\%\) PCK@0.10, while pair-level filtering yields a further improvement to about \(93\%\) PCK@0.10 on the retained pairs. We additionally provide a detailed per-semantic-label analysis in \autoref{Sup:per_label}, following the evaluation protocol of~\cite{wimmer2024back}, as well as a qualitative analysis of DINOv2-based semantic filtering and correspondence pruning in \autoref{sup:supp_semantic_analy}.


% \subsection{Ablations and Hyperparameter Sensitivity}
% \label{sec:ablation}
% We summarize the main findings of our ablation studies here, and defer detailed tables and extended quantitative analyses to the supplementary material.

% \noindent\textbf{Hyperparameters.}
% Performance is stable across a reasonable range of the CPA smoothness weight \(\lambda_{\mathrm{smooth}}\) and semantic threshold \(\tau_{\mathrm{sem}}\). We observe only minor variation across tested \(\lambda_{\mathrm{smooth}}\) values, while \(\tau_{\mathrm{sem}}=0.85\) gives the best average PCK@0.10, PCK@0.15, and PCK@0.20, with a favorable accuracy--coverage trade-off. We therefore use \(\lambda_{\mathrm{smooth}}=0.3\) and \(\tau_{\mathrm{sem}}=0.85\) in all experiments.

% \noindent\textbf{Registration model.}
% We compare continuous piecewise-affine (CPA) registration with a global affine transformation and thin-plate splines (TPS). Affine performs substantially worse at all PCK thresholds, while TPS is much closer to CPA. CPA nonetheless achieves the best accuracy overall and is approximately \(3\times\) faster in our setup, likely because its tetrahedral piecewise-affine parameterization is local, whereas TPS defines a global radial-basis warp.

% \noindent\textbf{Semantic filtering.}
% Semantic filtering consistently improves correspondence quality. When we disable both correspondence-level semantic pruning and pair-level semantic filtering, geometric registration alone achieves over 84\% PCK@0.10. Adding correspondence-level pruning improves performance to about 91\% PCK@0.10, and pair-level filtering yields a further gain to about 92\% PCK@0.10 on the retained pairs. We also provide a finer-grained per-label analysis in the supplementary material following~\cite{wimmer2024back}.

